\section{A maximum law for orthogonal sums}\label{sec:sum}
Let
\begin{equation}\label{eq:decomposition}
 V=V_0\oplus V_1\oplus\cdots\oplus V_m
\end{equation}
be an orthogonal decomposition, where $V_0=V^G$ and each $V_i$,
$i\ge1$, is a nontrivial real irreducible representation. Multiplicities
are allowed. Put $d_i=\dim V_i$, $D=\dim V$ and
\begin{equation}\label{eq:sigma}
 \sigma(U)=\max_{1\le i\le m}s(U|_{V_i}),
 \qquad \sigma(U)=0\quad\hbox{if }m=0.
\end{equation}
Equivalent irreducible orthogonal representations have isomorphic
orbits. Thus $\sigma(U)$ does not depend on the chosen decomposition
of an isotypic summand. In general it is \emph{not} the smallest
orbit dimension on the sphere of the whole reducible representation.

For the lower bound assume that the finite unit seeds span $V$ and
write
\begin{equation}\label{eq:seed-activation}
 a_i=\max_{x\in\mathcal X}\norm{\pi_i x}>0,
 \qquad a_* =\min_{1\le i\le m}a_i,
\end{equation}
where $\pi_i$ is orthogonal projection; set $a_*=1$ if $m=0$.
In a basis-adapted seed set
$\{\pm e_j\}$ all $a_i$ equal one. The query basis need not respect
the decomposition.

\begin{theorem}[Irreducible maximum principle]\label{thm:sum-width}
Suppose a law on the finite command alphabet has a full group norm
gap. Let the seeds span $V$, let
$0<\rho\le\min\{1/10,1/(8D)\}$, and, when $m>0$, let
$0\le\varepsilon<\rho a_*/(2\sqrt D)$. Then
\begin{equation}\label{eq:sum-width}
                   W_{N,\varepsilon}=\Theta(N^{\sigma(U)/2}).
\end{equation}
The constants may depend on the representation, command law, seed
set, signal and error, but not on $N$. If $m=0$, an exact constant-width
machine exists. The upper realization uses common command rows and
only one selected terminal numerical query.
\end{theorem}
\begin{proof}
First fix $i\ge1$ and a seed $x$ with
$\norm{\pi_ix}=a_i$. Track the unit-vector process
$Y_t^{(i)}=U_{a_t}\cdots U_{a_1}\pi_i x/a_i$ in $V_i$ and its
conditional mean given the \emph{actual whole-machine state}. Form
the ambient vector $d=(d_1,\ldots,d_D)$ from the terminal columns in
the proof. It has norm at most $\sqrt D$. The prescribed row means
and the uniform error give
\[
 \E\langle Y_N^{(i)},\pi_i d(S_N)\rangle
                \ge\rho a_i-2\sqrt D\varepsilon.
\]
Conditioning on $S_N$ bounds the left side above by
$\sqrt D\E\norm{\E[Y_N^{(i)}\mid S_N]}$.
Hence the terminal conditional amplitude is at least
$\rho a_i/\sqrt D-2\varepsilon>0$. The inherited entropy occupation
argument applies in the representation $V_i$ to these conditional
means, even though the hidden label also serves every other
constituent. Theorem~\ref{thm:uniform-orbits} supplies its exponent
$s(U|_{V_i})$. It follows that
$W_{N,\varepsilon}\ge c_iN^{s(U|_{V_i})/2}$.
Choose a constituent attaining \eqref{eq:sigma}.

For the upper bound include only nonzero summands and set $w_i=d_i/D$.
Choose a branch $i$ at initialization with probability $w_i$,
independently of the seed. The branch identifier is part of the state
label, not a free future input. In branch $i\ge1$ use the minimal-orbit
polytope of Theorem~\ref{thm:orbit-upper}. Its initial represented mean
is prescribed to be
\[
                       (\rho/w_i)\pi_i x.
\]
Its norm is at most $\rho D/d_i\le1/(8d_i)$, so the initialization
proof there applies also to this possibly subunit vector. Subsequent
rows are the common command rows from that construction. At a query
$j$ output the mean $\langle e_j,a_Nv_{i,s}\rangle$. The conditional
represented vector in the branch is
$(\rho/w_i)U_w\pi_ix$; multiplying by $w_i$ and summing gives exactly
$\rho U_wx$. All terminal means are bounded by $2/7$.

In the invariant branch, if present, use the fixed vectors
$\{\pm e:e\hbox{ in an orthonormal basis of }V_0\}$ scaled by
$1/4$, with identity transitions. Their convex hull contains a ball
of radius $1/(4\sqrt{d_0})$, which contains the required initial
vector of norm at most $1/(8d_0)$. Its contribution is exact with
$2d_0$ labels. The entire register is the disjoint union of these
branch labels. Its cardinality is at most
\[
 2d_0+\sum_{i=1}^m C_i N^{s(U|_{V_i})/2}
                   =O(N^{\sigma(U)/2}).
\]
No product of independently correct outputs is required. If $m=0$
the invariant construction proves the asserted constant bound. The
fixed two-label answer cut does not change the orders.
\end{proof}

The original full-group maximum principle is retained here in its
original sufficient range. Theorem~\ref{thm:intrinsic-width41} strengthens
it by using the active span, canonical activation and only one dominant
action gap; its proof of the current main theorem appears in
Section~\ref{sec:activation41}.

The maximum in \eqref{eq:sigma}, rather than a sum, is substantive.
For the direct product action of $\SO(3)\times\SU(3)$ on
$\R^3\oplus\Herm(3)$, the two least orbit dimensions are two and
four. With a full-gap product command law and the above signal
conventions the label order is $N^2$, not $N^3$; a generic product
orbit and a joint-output product task describe different quantities.
Adding a fixed subspace creates zero-dimensional whole-space orbits
but does not remove this $N^2$ requirement. The proof explains both
phenomena by paid branch labels and a constituent-wise converse.


\begin{proposition}[All separate cuts have constant positive minimum]
\label{prop:separate-width39}
For any orthogonal command family on $\R^D$, seeds
$\mathcal X=\{\pm e_1,\ldots,\pm e_D\}$, and $0<\rho\le1/D$,
the normalized residual row rank and the separate positive minimum
are both $D+1$ at every internal cut. The separate minimum permits
unrestricted memory at the other cuts. The assertion requires no
group gap and does not assert simultaneous attainment.
\end{proposition}
\begin{proof}
For a vector $z$, its remaining response on suffix $v$ and query $j$
has mean $\langle e_j,U_vz\rangle$. The resulting affine map from $z$
to the normalized response array is injective: fix one suffix and
vary the coordinate query; $U_v$ is invertible. At a fixed cut, fix
one preceding word $w$ and vary the seeds. The vectors
$\{\pm\rho U_we_j\}$ affinely span $\R^D$. Thus the normalized
residual arrays have row rank $D+1$, giving the positive lower bound.

A regular unit simplex in $\R^D$ contains the ball of radius $1/D$.
Every actual represented vector at that cut has norm $\rho$ and is
therefore a convex combination of its $D+1$ vertices. A vertex
$v_s$ defines a legal normalized continuation, since
$|\langle e_j,U_vv_s\rangle|\le1$ for every suffix and query.
These are $D+1$ positive generators of all cut rows. To embed the
factorization in a whole machine, retain the complete finite prefix
until the selected cut, sample its convex coordinates there, then
retain the chosen label together with the remaining command prefix.
There is no cap at those other cuts. The resulting finite machine
attains $D+1$ at the selected cut. This verifies the operational,
not merely amputated-matrix, meaning of the separate minimum.
\end{proof}
